\clearpage

\lehead[]{\normalfont\sffamily\hspace*{-2.00cm}\textcolor{white}{\colorbox{lightblue}{\parbox[c][0.70cm][b]{1.60cm}{
\makebox[1.60cm][r]{\thechapter}\\ \makebox[1.60cm][r]{ÜBUNG}}}}\hspace{0.17cm}\textcolor{lightblue}{\chaptertitle}}
\rohead[]{\textcolor{lightblue}{\chaptertitle}\normalfont\sffamily\hspace*{0.17cm}\textcolor{white}{\colorbox{lightblue}{\parbox[c][0.70cm][b]{1.60cm}{\thechapter\\
ÜBUNG}}}\hspace{-2.00cm}}
%\chead[]{}
\rehead[]{\textcolor{lightblue}{AvHG, Inf, My}}
\lohead[]{\textcolor{lightblue}{AvHG, Inf, My}}

\section{Erste Übung -- zweiter Teil}

Noch einmal zu unserem Problem aus der ersten Übung: Die Quadratwurzel einer
beliebigen positiven Zahl soll auf zwei Nachkommastellen genau bestimmt werden,
\emph{ohne}(!) die Wurzelfunktion des Taschenrechners zu benutzen.

\paragraph{Euer Auftrag:}
\begin{compactenum}[1.]
\item Setzt euch noch einmal in Zweier- oder Dreiergruppen zusammen.
\item Beschreibt euren Algorithmus diesmal im \textbf{Pseudocode} aus dem
vorherigen Kapitel -- also mithilfe von \textbf{Variablen} (z.\,B.\
\lstinline|zahl|, \lstinline|ergebnis|) und \textbf{Kontrollstrukturen}
(\lstinline|FALLS ... SONST|, \lstinline|WIEDERHOLE SOLANGE|).
\item Nutzt am Ende den Befehl \lstinline|PRINT|, um das fertige Ergebnis
auszugeben.
\end{compactenum}

\medskip
\noindent \textbf{Tipp:} Falls ihr beim Aufschreiben merkt, dass sich euer
Lösungsweg aus dem ersten Übungsteil nur umständlich in Schleifen und
Bedingungen fassen lässt, dürft ihr euren Algorithmus gerne anpassen oder
vereinfachen!

\medskip
\noindent \textbf{Zeit:} 15 Minuten \\
\textbf{Sicherung:} Haltet euren Pseudocode so sauber fest, dass eine andere
Gruppe ihn Schritt für Schritt am Taschenrechner nachvollziehen („durchspielen“)
kann.
